\section{职业与研究：在真实系统中制造 novel constraints}

学生问 Rauch 是否后悔过早进入工业界而没有走传统研究路线。他的回答不是“工业优于学术”，而是自己在公司里参与 scheduler、placement、load balancer、metadata store 等研究型工作。真实产品可以提供高规模、强约束和即时反馈；但他也承认过多 DIY 会杀死 startup。

\begin{knowledgebox}{Applied research 的三个必要条件}
\begin{enumerate}
  \item 问题包含现有方法不能满足的明确 constraint；
  \item 团队能构建实验、测量和失败反馈，而不是只写生产补丁；
  \item 研究结果能回到产品指标，同时保留可复用知识和系统边界。
\end{enumerate}
\end{knowledgebox}

\begin{warningbox}{“自己写一遍”不是默认教育路线}
重写 React、Kubernetes 或 scheduler 可以建立深层理解，也可能浪费公司 runway。应选择与核心差异化直接相关、现成方案无法满足约束的层；其余部分优先复用，并通过源码阅读、实验和故障演练学习。
\end{warningbox}

\teachervoice{课堂真正有张力的地方是两句话同时成立：too much DIY can kill you；unreasonable standards can create novel research. 成熟判断不在“全部买”或“全部造”，而在知道哪个约束值得用组织生命去解决。}

\subsection{本章小结}

工业系统可以成为研究环境，但前提是问题约束、实验反馈和产品价值清晰。Build vs. Buy 同样适用于个人学习：深挖关键层，复用非关键层。

